\documentclass[12pt,letterpaper]{article}
\usepackage[T1]{fontenc}
\usepackage[letterpaper,left=1in,right=1in,top=.94in,bottom=.86in,
  headheight=23pt,headsep=15pt,footskip=29pt]{geometry}
\usepackage{amsmath,amssymb,mathtools,newpxtext,newpxmath}
\usepackage[semibold]{sourcesans}
\usepackage[varqu]{inconsolata}
\usepackage{microtype,xcolor,graphicx,booktabs,tabularx,array,enumitem,fancyhdr}
\usepackage[most]{tcolorbox}
\usepackage{hyperref}
\definecolor{ink}{HTML}{19303A}
\definecolor{sea}{HTML}{00675E}
\definecolor{clay}{HTML}{A24725}
\definecolor{muted}{HTML}{4C636C}
\definecolor{rulegray}{HTML}{CAD4D4}
\definecolor{wash}{HTML}{EAF4F0}
\hypersetup{colorlinks=true,linkcolor=sea,urlcolor=sea,
  pdftitle={A prime sum and its zeros},
  pdfsubject={An illustrated companion to Levent Alpöge's proposed proof of Dirichlet GRH},
  pdfauthor={Codex, for the Dirichlet GRH sketch repository},bookmarksnumbered=false}
\urlstyle{same}
\color{ink}
\setlength{\parindent}{0pt}
\setlength{\parskip}{5pt plus .4pt minus .5pt}
\setlength{\abovedisplayskip}{6pt plus 1pt minus 2pt}
\setlength{\belowdisplayskip}{6pt plus 1pt minus 2pt}
\setlength{\abovedisplayshortskip}{4pt}
\setlength{\belowdisplayshortskip}{4pt}
\setlist{nosep,leftmargin=1.6em}
\setlist[enumerate]{itemsep=3pt}
\pagestyle{fancy}\fancyhf{}
\fancyhead[L]{\sffamily\footnotesize\color{muted}A prime sum and its zeros}
\newcommand{\partnamehere}{Reader's guide}
\fancyhead[R]{\sffamily\footnotesize\color{muted}\partnamehere}
\fancyfoot[L]{\sffamily\footnotesize\color{muted}Dirichlet GRH\enspace/\enspace An illustrated companion}
\fancyfoot[R]{\sffamily\bfseries\small\color{sea}\thepage}
\renewcommand{\headrulewidth}{.4pt}
\renewcommand{\headrule}{\hbox to\headwidth{\color{rulegray}\leaders\hrule height \headrulewidth\hfill}}
\fancypagestyle{opening}{\fancyhead[L]{\sffamily\footnotesize\color{muted}Dirichlet GRH sketch repository}
  \fancyhead[R]{\sffamily\footnotesize\color{muted}20 November 2026}}
\newcommand{\sheet}[3]{\clearpage\renewcommand{\partnamehere}{#1}\phantomsection
  \pdfbookmark[1]{#2}{bm-#3}\label{#3}
  {\sffamily\bfseries\fontsize{23}{26}\selectfont #2\par}\vspace{2pt}}
\newcommand{\subt}[1]{\par\addvspace{6pt}{\sffamily\bfseries\large\color{sea}#1\par}\nobreak\vspace{-2pt}}
\newcommand{\smallhead}[1]{{\sffamily\bfseries\color{sea}#1}}
\newcommand{\sources}[1]{\par\vfill\noindent{\color{rulegray}\rule{\linewidth}{.45pt}}\par
  \vspace{-4pt}{\footnotesize\sffamily\color{muted}#1\par}}
\newcommand{\mirror}{https://github.com/TomasOrtega/RH/blob/a6ac91a412dcbaab31801355ec5707024a71bf65/extracted/lean/}
\newcommand{\src}[3]{\href{\mirror src/#1.lean\#L#2}{#3}}
\newcommand{\pub}[2]{\href{https://alpo.ge/remark.pdf\#page=#1}{#2}}
\newcommand{\aaDir}{Development/AutomorphicArithmetic/}
\newcommand{\chDir}{Development/AutomorphicArithmetic/CharacterBoxes/}
\newcommand{\dcDir}{Development/AutomorphicArithmetic/DedekindLFunctions/}
\newcommand{\laDir}{Development/AutomorphicArithmetic/DedekindLFunctions/LandauTheorems/}
\newcommand{\prDir}{Development/AutomorphicArithmetic/PrimitiveCosets/}
\newcommand{\frDir}{Formalism/SharpRestriction/RamificationFrobenius/TotientGaloisClasses/}
\newcommand{\sizeDir}{Formalism/SizeWindowBrackets/}
\newcommand{\growthDir}{Formalism/SharpRestriction/GrowthBounds/}
\newcommand{\thirdDir}{Development/AutomorphicArithmetic/ThetaCorrespondence/ThetaPairLifts/ThirdPlaceSeries/}
\newcommand{\repo}{https://github.com/robbiemu/dirichlet-grh-sketch/}
\newcommand{\sep}{\enspace{\color{rulegray}\textbar}\enspace}
\newtcolorbox{statement}{enhanced,sharp corners,boxrule=0pt,colback=wash,
  colframe=wash,borderline west={2pt}{0pt}{sea},left=9pt,right=9pt,top=6pt,
  bottom=6pt,before skip=5pt,after skip=5pt}
\newcounter{illustration}
\newcommand{\ill}[3][.99]{\par\addvspace{5pt}\noindent
  \begin{minipage}{\linewidth}\centering
    \includegraphics[width=#1\linewidth]{figures/#2.pdf}\par\vspace{3pt}
    \raggedright\refstepcounter{illustration}
    {\footnotesize\sffamily\color{muted}\textbf{\color{sea}Figure \theillustration.}\enspace #3\par}
  \end{minipage}\par}
\newcommand{\R}{\mathbb R}\newcommand{\C}{\mathbb C}\newcommand{\Z}{\mathbb Z}
\newcommand{\Q}{\mathbb Q}\newcommand{\F}{\mathbb F}\newcommand{\one}{\mathbf1}
\newcommand{\re}{\operatorname{Re}}\newcommand{\Res}{\operatorname*{Res}}
\newcommand{\W}{\mathrm{W}}\newcommand{\e}{\varepsilon}
\newcommand{\Lder}{\frac{L'}{L}}
\newcommand{\pagego}[2]{\hyperref[#1]{#2}\hfill\pageref{#1}}
\begin{document}
\thispagestyle{opening}\pdfbookmark[1]{A prime sum and its zeros}{opening}
{\sffamily\bfseries\fontsize{38}{42}\selectfont A prime sum and its zeros\par}
{\sffamily\fontsize{16}{20}\selectfont\color{sea}
An illustrated companion to Levent Alpöge's\\proposed proof of Dirichlet GRH\par}
\vspace{7pt}
The generalized Riemann hypothesis for Dirichlet $L$-functions (Dirichlet
GRH) says that for every character $\chi$, every zero $\rho$ of $L(s,\chi)$
in the strip $0<\re\rho<1$ has $\re\rho=\tfrac12$. At conductor one this
is the Riemann hypothesis for the zeta function.

Can a sum over ordinary prime powers detect a zero off that line?
In this argument it would detect one by how often the sum rises above a threshold.
A second construction, using automorphic forms, is meant to keep the same sum
below the threshold across some very long intervals. The contradiction is about
\emph{where} a large value must occur.

\ill{windows}{The horizontal coordinate is $u=\log x$. The dot marks a possible
  location of a violation. Both rows concern the same sum and threshold.}

\subt{How to read the argument}
The first part constructs the sum and explains why a zero has consequences on
every sufficiently late interval. The second asks what arithmetic information
could produce even one good interval beyond each prescribed starting point.
A reader familiar with calculus, linear algebra and finite groups can begin at the
first line. Complex analysis first enters in the discussion of the transform.

{\sffamily\small
\begin{tabularx}{\linewidth}{@{}p{1.52in}X@{}}
\toprule
Prime sums & \pagego{characters}{Characters, ramps and the window estimate}\\
Complex analysis & \pagego{transform}{How a zero survives smoothing and averaging}\\
Arithmetic & \pagego{finite}{From a finite Fourier calculation to iterated lifts}\\
Assembly & \pagego{frobenius}{How the estimates reach the prime sum}\\
Verification & \pagego{lean}{How to read and check the released statements}\\
\bottomrule
\end{tabularx}}

\sources{This independent exposition expands the repository's
\href{\repo blob/main/proof-sketch.pdf}{four-page sketch}. Alpöge describes the proposal and its
formalization in his \pub{1}{announcement}. The calculations and models below can be
reproduced. The uniform arithmetic construction and an end-to-end Lean kernel check
remain outside the results established here.}

\sheet{I\enspace Prime sums}{Characters and residue classes}{characters}
Among the odd primes, $5$ and $13$ are $1$ modulo $4$, while $3$ and $7$ are
$3$ modulo $4$. Give the first class weight $+1$ and the second weight $-1$.
The resulting signed sum records the difference between the classes. Its
weights multiply because $3\cdot3\equiv1\bmod4$.

\subt{A finite Fourier coordinate}
For an integer $q\ge1$, a Dirichlet character $\chi$ is a homomorphism from
the units $(\Z/q\Z)^\times$ to $\C^\times$, extended periodically to the
integers and set to zero on nonunits. The values on units are roots of unity.
The \emph{principal} character has value $1$ on every unit. The nonprincipal
character $\chi_4$ above and the principal $\chi_{0,4}$ have the table
\[
\begin{array}{c|rrrr}
n\bmod4&0&1&2&3\\\hline
\chi_{0,4}(n)&0&1&0&1\\
\chi_4(n)&0&1&0&-1
\end{array}
\]
For instance, $\one_{n\equiv1\,(4)}=(\chi_{0,4}(n)+\chi_4(n))/2$, including
at even $n$. Character sums let us recover a count for an individual class.

\subt{Why a product over primes knows about their powers}
Write $s=\sigma+it$ and $n^{-s}=\exp(-s\log n)$, using natural logarithms.
For $\sigma>1$, absolute convergence and unique prime factorization give
\[
L(s,\chi)=\sum_{n\ge1}\frac{\chi(n)}{n^s}
=\prod_p (1-\chi(p)p^{-s})^{-1}.
\]
Expand the logarithm of each factor and differentiate term by term. Cancelling
the $r$ from the derivative of $p^{-rs}$ against the $1/r$ in the logarithm gives
\[
 -\Lder(s,\chi)=\sum_p\sum_{r\ge1}\frac{\chi(p)^r\log p}{p^{rs}}
 =\sum_{n\ge1}\frac{\Lambda(n)\chi(n)}{n^s}.
\]
Here $\Lambda(p^r)=\log p$ for $r\ge1$ and $\Lambda(n)=0$ otherwise, including
at $n=1$. Thus $9$ carries weight $\log3$, and $6$ carries weight zero.

\subt{What happens beyond convergence}
Analytic continuation gives meaning to $L(s,\chi)$ in the strip $0<\sigma<1$.
The conductor is the least modulus from which the character is induced. A
character at its conductor is \emph{primitive}. The principal character of
conductor one gives $\zeta(s)$, with a simple pole at $1$. A nonprincipal
Dirichlet $L$-function is regular and nonzero at $1$. These classical facts
determine what must be subtracted from the prime sum on the next page.

\sources{For the character expansion, continuation and functional equation see the
\href{https://dlmf.nist.gov/25.15}{NIST Digital Library of Mathematical Functions, §25.15}.
For the exact strip statement used in the release see \src{Statements}{10}{Statements.lean}.}

\sheet{I\enspace Prime sums}{Giving each prime power a ramp}{ramp}
A sharp cutoff includes a new term in a sum the instant $x$ reaches $n$.
Instead, give that term a weight that starts at zero and increases linearly.
For $x\ge1$ the resulting continuous, piecewise linear function is
\[
 R_\chi(x)=\sum_{n\le x}\Lambda(n)\chi(n)(x/n-1).
\]
Its slope changes by $\Lambda(n)\chi(n)/n$ at $n$. This choice of shape will
make its integral transform a particularly simple multiple of $-L'/L$.

\subt{Two sums that can be checked by hand}
At $x=4$ the conductor-one ramp has received contributions from $2$ and $3$.
The term for $4$ is still zero, so $R_{\one}(4)=\log2+\tfrac13\log3$.
For $\chi_4$ at $x=10$, the contributing prime powers are $3,5,7,9$. Thus
\[
 R_{\chi_4}(10)
 =-\tfrac73\log3+\log5-\tfrac37\log7+\tfrac19\log3 .
\]
The $9$ term has positive sign because $\chi_4(3)^2=1$.

\ill{prime-ramp}{Direct sums at quarter-integer spacings. At left, interpolation is
exact between adjacent integers. At right, straight segments join samples of a
curved function. The panels have different axes.}

\subt{The fluctuation we want to bound}
Let $\gamma=\lim_{n\to\infty}(\sum_{j=1}^n1/j-\log n)$. For the conductor-one
character $\one$ and for nonprincipal $\chi$, respectively, put
\[
\begin{aligned}
 G_{\one}(x)&=R_{\one}(x)-x\log x+(1+\gamma)x,\\
 G_\chi(x)&=R_\chi(x)+x\,L'(1,\chi)/L(1,\chi).
\end{aligned}
\]
For example $G_{\one}(4)\simeq1.82304$. Although $R_{\one}$ has grown past
$12$ by $x=12$, its centered value is of order one in this example. The
subtractions reflect an exact pole cancellation (p.~\pageref{transform}).
The graph covers a finite range. An upper bound on arbitrarily late intervals
demands an argument at unbounded $x$.

\sources{The released definitions are \src{\dcDir CharFamily}{10}{ramp and Gchi}.
The accompanying \texttt{figures/figure\_data.py} generates the numerical table
by sieving for prime powers.}

\sheet{I\enspace Prime sums}{One fixed sum, an entire window}{window}
Fix a primitive target $\psi$ of order $h$, the least positive power
giving the principal character. Take the inducing primitive characters
of all powers of $\psi$ at its modulus. The resulting set
$F_\psi$ contains $\psi$, its inverse and $\one$. The release measures
\[
 H_\psi(x)=G_{\one}(x)
       +\frac1{3h+1}\sum_{\chi\in F_\psi\setminus\{\one\}}\re G_\chi(x).
\]
Taking real parts gives a sum we can order. The positive weights preserve a
target pole and arise from averaging over residue classes (p.~\pageref{frobenius}).

\subt{An example with a nonreal character}
The units modulo $5$ are generated by $2$. Choose $\psi(2)=i$, so $\psi$
has values $1,i,-1,-i$ on $1,2,4,3$ in that order. Its square is the
quadratic character, its cube is $\bar\psi$, and its fourth power induces
$\one$ at conductor one. Here $h=4$ and the weight is $1/13$ for each of
the three nonprincipal characters. For $\chi_4$ the weight is $1/7$.
The family and weights stay fixed as the window moves.

\begin{statement}
\smallhead{The window estimate $(\W)$}\quad
For every primitive $\psi$, every real $\theta$ with
$\tfrac12<\theta<1$, and every $Y_0\ge2$, there exists a real $Y\ge Y_0$
such that
\[
           H_\psi(x)\le x^\theta \qquad\text{for every real }x\in[Y,Y^{3/2}].
\]
\end{statement}
At $Y=100$ this covers every point of $[100,1000]$. The interval may be
chosen after $\psi,\theta,Y_0$.

\subt{What a contradiction would look like}
Suppose $L(\rho,\psi)=0$ and $\tfrac12<\beta=\re\rho<1$. Fix, for example,
$\theta=(\beta+\tfrac12)/2$. The analytic criterion implies that there is a
$Y_*$ such that \emph{for every} $Y\ge Y_*$ there is \emph{some}
$x\in[Y,Y^{3/2}]$ with $H_\psi(x)>x^\theta$. Apply $(\W)$ with
$Y_0=\max(2,Y_*)$. The violating point lies inside the very interval on
which the arithmetic upper bound holds.

\ill[.84]{strip}{Critical strip with hypothetical $\rho=3/4+i/2$ and $\theta=5/8$.
The functional equation pairs $\rho$ for $\psi$ with $1-\rho$ for $\bar\psi$
(p.~\pageref{assembly}). The coordinates are chosen to show the symmetry.}

\sources{The cyclic family is built in \src{\frDir ExplicitFormulaErrorOfHostConditions}{105}{Level.ofSub and Level.cyc}.
See also its \src{\prDir LevelFamily}{41}{conversion to a positive family}.
The exponent $3/2$ is chosen in \src{\chDir EveryCharacter}{8}{plugWin}.}

\sheet{II\enspace The analytic argument}{Turning a ramp into a pole}{transform}
The Mellin transform of a function on $[1,\infty)$ is
$\widehat f(s)=\int_1^\infty f(x)x^{-s-1}\,dx$ wherever the integral converges.
It converts powers of $x$ into simple fractions. For instance,
$\widehat{x^\alpha}(s)=1/(s-\alpha)$ when $\re s>\re\alpha$.
This is a way to read an exponent from a singularity. More generally,
near a simple pole at $\rho$ a function has the form $c/(s-\rho)+h(s)$,
where $h$ has a convergent local power series. The number $c$ is the
\emph{residue}. A meromorphic continuation extends the function analytically
while allowing isolated poles.

\subt{Integrate one summand}
For $\re s>1$, integrating the ramp that starts at $n$ gives
\[
\int_n^\infty(x/n-1)x^{-s-1}\,dx
  =n^{-s}\left(\frac1{s-1}-\frac1s\right)
  =\frac{n^{-s}}{s(s-1)}.
\]
The sum of the absolute integrals converges there because
$|\chi(n)|\le1$ and $\sum\Lambda(n)n^{-\sigma}<\infty$ for $\sigma>1$.
We can therefore interchange integration and summation and obtain
\[
  \widehat R_\chi(s)=-\frac{1}{s(s-1)}\Lder(s,\chi).
\]
The right-hand side gives a meromorphic continuation. Integrals will only be
differentiated in a half-plane of convergence.

\subt{Cancel the singularity at one}
Put $z=s-1$. From $\zeta(1+z)=1/z+\gamma+O(z)$ we get
$-\zeta'/\zeta=1/z-\gamma+O(z)$. Multiplication by
$1/[s(s-1)]=1/z-1+O(z)$ gives the two singular terms
$z^{-2}-(1+\gamma)z^{-1}$. The transforms of $x\log x$ and $x$ are $z^{-2}$
and $z^{-1}$, respectively. This gives exactly the two corrections in
$G_{\one}$. For nonprincipal $\chi$, regularity and nonvanishing at one leave
just $-(L'/L)(1,\chi)/z$ to cancel. This explains the plus sign in $G_\chi$.

\subt{Keep the singularity at a zero}
At a strip zero $\rho$ of multiplicity $r$, write
$L(s,\chi)=(s-\rho)^r h(s)$ with $h$ holomorphic and $h(\rho)\ne0$.
Then $L'/L=r/(s-\rho)+h'/h$. Both centering terms are holomorphic at $\rho$,
so the continued transform has a simple pole with residue
\begin{statement}
\[
             \Res_{s=\rho}\widehat G_\chi(s)
                 =-\frac{r}{\rho(\rho-1)}.
\]
\end{statement}
For a real $0<\rho<1$ this number is positive. For a complex $\rho$ the
phase depends on $\rho$, and all characters sharing that zero have the
same phase. This fact is what allows the fixed family to keep the zero visible.

\sources{Primary calculations in the release include \src{\laDir MellinRamp}{10}{the ramp integral},
\src{\laDir mellinGchiModOne}{115}{conductor-one centering},
\src{\laDir mellinGchiOfNeOne}{116}{nonprincipal centering} and
\src{\laDir poleData}{137}{the local pole expansion}.}

\sheet{II\enspace The analytic argument}{Why the average still detects a zero}{survival}
Let $F$ be any fixed finite family with positive real weights $a_\chi$.
Admit a principal character only at conductor one. Write
$H_F=\sum_{\chi\in F}a_\chi\re G_\chi$ and $\Phi=\widehat H_F$. Since
$G_{\bar\chi}(x)=\overline{G_\chi(x)}$ for real $x$, the transform of
$\re G_\chi$ is $(\widehat G_\chi(s)+\widehat G_{\bar\chi}(s))/2$. The
two transforms here are evaluated at the \emph{same} complex argument $s$.

Let $r_\chi(\rho)$ be the order of vanishing, with value zero when
$L(\rho,\chi)\ne0$. At a zero of a family member the residue in $\Phi$ is
\[
  \frac{-1}{\rho(\rho-1)}
    \sum_{\chi\in F}\frac{a_\chi}{2}
    \bigl(r_\chi(\rho)+r_{\bar\chi}(\rho)\bigr).
\]
The sum on the right is a positive real number. Thus $\rho$ persists in the
transform and any real residue in the open strip is positive.

\subt{An oscillation we can solve exactly}
Take the model $f(x)=x^{3/4}\cos(2\pi\log x)$. Its Mellin transform has
two poles, at $3/4\pm2\pi i$, since $f$ is the average of the two powers
$x^{3/4\pm2\pi i}$. To compare $f$ with $x^{3/5}$, set $u=\log x$.
The ratio is $e^{3u/20}\cos(2\pi u)$, and at each positive integer $u$ it
exceeds $1$. Once $U\ge2$, the interval $[U,3U/2]$ contains an integer.
Every such interval therefore supplies a violation.

\ill{oscillation}{A two-pole model. Shading marks $[2,3]$ and $[4,6]$ in log coordinates.
The curve plots $e^{3u/20}\cos(2\pi u)$. The dots mark integer values of $u$ inside the windows.}

For an actual family the continued transform can have many poles. Their
oscillations can interfere, and positivity of the residue coefficient at a
given pole says nothing immediate about the sign of $H_F$ at a prescribed
$x$. To obtain a statement on every late window, one needs both localization
in $u$ and separation of poles in the $s$-plane.

\sources{The source combines the transforms for a character and its inverse in
\src{\laDir poleData}{248}{poleData}. The equality of the conjugated prime sums is
recorded in \src{\laDir GchiConj}{38}{GchiConj}.}

\sheet{II\enspace The analytic argument}{Localizing a one-sided bound}{kernel}
Put $g(u)=H_F(e^u)$ for $u\ge0$. Then $\Phi(s)=\int_0^\infty g(u)e^{-su}\,du$.
Summing $\log n/n$ up to $x$ gives $O((1+\log x)^2)$, so the ramps and their
centerings give $|g(u)|\le B e^{2u}$ for some $B$ depending on $F$.
For $\sigma>2$, differentiation under the integral gives
\[
 A_k(\sigma+it):=\frac{(-1)^k}{k!}\Phi^{(k)}(\sigma+it)
 =\int_0^\infty g(u)e^{-itu}\,w_k(u)\,du,\qquad
 w_k(u)=\frac{u^k e^{-\sigma u}}{k!}.
\]
The triangle inequality and the reality of $g$ give the one-sided test
\[
 |A_k(\sigma+it)|+\re A_k(\sigma)
       \le \int_0^\infty (|g(u)|+g(u))w_k(u)\,du.
\]
The integrand on the right sees only positive values of $g$. On a good window
$[U,3U/2]$, where $g(u)\le C e^{\theta u}$ with $C>0$, its integral is at
most $2C/(\sigma-\theta)^{k+1}$, by integration of $u^k e^{-(\sigma-\theta)u}/k!$.

\ill{kernel}{The weight divided by its maximum, at $u/U=5/4$, for two choices used
in the tail estimate. The shaded window is fixed in these rescaled coordinates.}

\subt{A tail calculation for this window length}
The release's choices specialize to $\sigma=336$, $k=\lceil420U\rceil$ for
$U\ge1$. Outside the window, multiplying the weight by the crude growth bound
leaves $u^ke^{-334u}/k!$. On $u>3U/2$, majorize it by multiplying by
$e^{14(u-3U/2)}$. On $u<U$, use $e^{14(U-u)}$. Extending the integrals of
the resulting positive functions over all $u>0$ bounds their sum by
\[
      \frac{e^{-21U}}{320^{k+1}}+\frac{e^{14U}}{348^{k+1}}
           \ \ll\ 336^{-k}.
\]
Indeed $420\log(336/320)<21$ and $420\log(348/336)>14$.
Consequently $|A_k(\sigma+it)|+\re A_k(\sigma)\le C_1(\sigma-\theta)^{-k}$
uniformly in $t$ along arbitrarily late good windows. Here $\ll$ hides a constant
independent of $U,k,t$. On this sequence $k\to\infty$.

\sources{The numerical choices come from \src{\laDir landauTails}{80}{landau\_tails}
at $\delta=1/2$. Compare the \src{\laDir landauWindowIneq}{27}{one-sided integral inequality}
and \src{\laDir landauWindow}{40}{its use for an integral majorant}.}

\sheet{II\enspace The analytic argument}{A pole defeats the window bound}{separation}
Why differentiate to high order? At a center $z$ the term $c/(s-\rho)$
contributes exactly $c/(z-\rho)^{k+1}$ to $A_k(z)$. A holomorphic remainder
on a closed disk of radius $r$ contributes at most $M/r^k$ by Cauchy's
estimate. If $d=|z-\rho|<r$, the pole eventually dominates. The extra
$\re A_k(\sigma)$ in our one-sided test also needs control.

\subt{Choose the two centers together}
Set $E=\sigma-\theta$ and assume there is a pole with real part $>\theta$.
The poles form a locally finite set in $\re s>0$, and lie to the left of
$\sigma$. Let $D$ be the distance from the real center $\sigma$ to the nearest
pole in its closed disk of radius $E$, taking $D=E$ if the disk is empty.
There are two possibilities.

\begin{enumerate}
\item If some pole $\rho_1$ satisfies $\sigma-\re\rho_1<D$, move the other
center to $\sigma+i\operatorname{Im}\rho_1$ and perturb its height slightly.
Its nearest pole then has distance $d<D$ and can be made unique. Only finitely
many poles can compete in a small neighborhood of that center. Each pair
ties at at most one height, since all poles lie to the left of the vertical line.
Choose $d<r<D$ also smaller than the distance to any other pole there.
The shifted center detects its single pole, while the disk of radius $r$
at the real center is pole-free.
\item Otherwise every pole has horizontal distance at least $D$. The assumed
pole to the right of $\theta$ implies $D<E$. A nearest pole at distance $D$
must therefore have zero vertical displacement. It is the unique real point
$\beta=\sigma-D$. Isolate it at the real center with $d=D<r<E$.
Its positive residue makes $\re A_k(\sigma)$ positive to leading order.
\end{enumerate}
In both cases the left side of the test on p.~\pageref{kernel} is bounded below
by $c_0d^{-k}-M'r^{-k}$ with $c_0>0$. An upper bound $C'E^{-k}$ along
unbounded $k$ would give
\[
        c_0\le M'(d/r)^k+C'(d/E)^k\longrightarrow0 .
\]
This proves the contradiction for the pointwise window bound used here.

\begin{statement}
\smallhead{Window criterion.}\quad For a fixed finite positive family as on
p.~\pageref{survival}, $0<\theta<1$, a fixed $\lambda>1$ and a fixed $C>0$,
an upper bound $H_F(x)\le Cx^\theta$ throughout arbitrarily late
$[Y,Y^\lambda]$ excludes every strip zero of a member with $\re\rho>\theta$.
\end{statement}
For general $\lambda$, the release chooses $\sigma$ to fit $\lambda-1$.
Its integral-majorant formulation allows a factor $e^{\eta U}$ for every
$\eta>0$. Taking $\eta<\log(E/d)$ absorbs that factor since $k\ge U$.

\sources{Follow \src{\laDir LandauDichotomy}{24}{the two cases},
\src{\laDir LandauNearest}{36}{choice of a unique nearest pole},
\src{\laDir landauWindow}{83}{pole comparison} and
\src{\laDir ZeroFreeOfRpowWindow}{11}{the pointwise corollary}.}

\sheet{III\enspace The arithmetic construction}{A Fourier calculation over a finite field}{finite}
Why introduce a chain? The announcement starts with an identity in which
neighboring variables interact. Over a field $\F_q$ of $q$ elements choose a
nontrivial additive character $e_q$, so $\sum_a e_q(az)$ equals $q$ when
$z=0$ and zero otherwise. For prime $q$ take $e_q(z)=e^{2\pi iz/q}$.
For a function $t:\F_q\to\C$ let $M=\sum_x t(x)$. On $k\ge2$ variables set
\[
 Q(\mathbf x)=\sum_{j=1}^{k-1}x_jx_{j+1},\qquad
 T(\mathbf x)=\prod_{j=1}^{k}t(x_j),\qquad N=\sum_{Q(\mathbf x)=0}T(\mathbf x).
\]
Summing over the frequency $a$ imposes the equation $Q=0$. Separating the
frequency $a=0$ instead makes the variables independent. The same sum gives
\[
             qN=M^k+R,\qquad
 R=\sum_{a\ne0}\sum_{\mathbf x}T(\mathbf x)e_q(aQ(\mathbf x)).
\]

\ill[.95]{path}{An exact example over $\F_5$. The right panel lists $Q$ modulo $5$ for
$t=\one_{\{1,4\}}$ and $k=3$. Four triples give zero, hence $N=4$,
$M=2$ and $R=5\cdot4-2^3=12$.}

\subt{Why the edges of a path matter}
For $a\ne0$ let $U_{xy}=q^{-1/2}e_q(axy)$. Orthogonality gives
$(UU^*)_{xz}=q^{-1}\sum_y e_q(a(x-z)y)=\one_{x=z}$, so $U$ preserves
Euclidean norm. Let $D=\operatorname{diag}(t)$, and let $v$ be the column of
$t$-values. The inner sum is $q^{(k-1)/2}v^{\mathsf T}U(DU)^{k-2}v$.
If $|t|\le1$, then $\|D\|\le1$ and $\|v\|\le\sqrt q$. Cauchy--Schwarz
therefore bounds each nonzero-frequency sum by $q^{(k+1)/2}$.

This calculation controls the remainder. To deduce cancellation in $M$
one also needs information about $N$. For instance $t\equiv1$ gives $M=q$
while satisfying exactly the same norm assumptions. The note describes a
geometric sign for its special trace functions. The path identity and the
eight-triple example establish no such sign. Over the integers, finding a
replacement for this geometric input is the reason the construction turns
to periods and divisibility.

\sources{The finite-field identity, the path and the proposed role of geometry appear
in the \pub{3}{announcement, p.~3, footnote 3}. The example evaluates its defining sum directly.}

\sheet{III\enspace The arithmetic construction}{When a zero is worth more than one}{multiplicity}
In a polynomial point count, an auxiliary polynomial can be made to vanish
at every point being counted. Suppose $P$ is a nonzero polynomial in one
variable of degree $d$, with $N$ distinct roots each of multiplicity at
least $m$. The corresponding $m$th powers of the linear factors divide $P$,
so $mN\le d$. Over $\F_5$, $P(X)=(X^2-1)^3$ counts $1$ and $4$ three
times each. Its degree $6$ gives equality. The requirement $P\ne0$ matters.

In Stepanov's method on a curve, one bounds the degree of a pole divisor of
an auxiliary function. A nonzero rational function has zero and pole
divisors of equal degree, giving the same kind of budget. Forcing larger
orders is worth doing only if the cost of constructing the function grows
slowly enough. This principle survives the passage to number fields.

\subt{The norm measures divisibility}
A number field $K$ is a finite extension of $\Q$. Its algebraic integers
are the elements satisfying a monic polynomial with integer coefficients.
Take a nonzero algebraic integer $a$.
Factor $(a)$ into prime ideals $\mathfrak p$, writing $v_{\mathfrak p}(a)$
for each exponent and $N\mathfrak p$ for the size of the residue field.
Two evaluations of the absolute norm give
\[
 \sum_{\mathfrak p}v_{\mathfrak p}(a)\log N\mathfrak p
       =\log|N_{K/\Q}(a)|=\sum_{\sigma:K\hookrightarrow\C}\log|\sigma(a)|.
\]
The last sum counts both members of a complex-conjugate pair. In $\Z[i]$
take $a=3+4i=(2+i)^2$. The ideal $(2+i)$ has norm $5$. The left side is
$2\log5$, and each of the two complex embeddings gives absolute value $5$.

\subt{Where repeated divisibility might come from}
The announcement gives a local prototype on a CM elliptic curve. Here CM
means that the curve has extra endomorphisms from an imaginary quadratic
field. At a good inert prime $p$, the note reports $D^p\equiv0\bmod p$ for
the invariant derivation $D$ in its integral setting. Inert means that $p$
stays prime in the quadratic field. Applied to an integral function, such a
congruence makes a string of derivatives acquire a factor of $p$.

The global proposal uses a determinant of periods. Linear relations modulo a
prime make a determinant divisible there. High multiplicity calls for
quantitative control of \emph{how many} relations, at which primes and to
what orders. The determinant also has to survive over the number field so that
its size can be compared with its divisibility. Alpöge reports first reaching
multiplicity two on $\mathrm O(2,n)$, the group preserving a real quadratic
form of signature $(2,n)$. The next construction tries to make this
multiplicity grow with the number of lifts.

\sources{Alpöge describes the differential congruence and the first determinant in
the \pub{3}{announcement, items 1--4}, and the iterated version in
\pub{4}{items 7--10}. The polynomial and Gaussian-integer computations are models
with their own exact conclusions.}

\sheet{III\enspace The arithmetic construction}{What travels through a theta lift}{lifts}
An automorphic form is a function on a group with arithmetic symmetries and
regularity conditions. A classical example lives on the upper half-plane
and transforms with a prescribed weight under integral matrices. Here the
groups are unitary groups $U(V)$ preserving a Hermitian form on a vector
space over a CM field, a totally imaginary quadratic extension of a totally
real number field. Total definiteness makes the unitary groups compact
at the real places. The finite primes still carry arithmetic information.
A \emph{level} specifies invariance under subgroups at those finite primes.

\subt{A kernel that can see two groups}
The classical series $\sum_{(a,b)\in\Z^2}e^{-\pi t(a^2+b^2)}$, $t>0$,
records the number of lattice vectors on each circle. For instance the
coefficient of $e^{-5\pi t}$ is eight, from $(\pm1,\pm2)$ and their swaps.
A theta kernel also sums over rational vectors, now with actions of two
groups on the underlying data. If $g$ and $h$ are variables on the old and
new groups, an integral transform has the schematic shape
\[
        f_{j+1}(h)=\int_{[U(V_j)]} f_j(g)\,\Theta_{j+1}(g,h)\,dg .
\]
Brackets denote the quotient by rational points in the adelic group, which
combines all completions of the number field. The kernel, pairings, measures
and choices of test vectors must be specified to evaluate actual periods.

\ill{lift}{The proposed chain contains $m$ transforms. The diagram records the
groups and the order of integration. Nonvanishing and integrality at the
successive stages are mathematical requirements.}

\subt{From a function to a number}
A period pairs a form with a character or another form and integrates along a
subgroup. It probes how a representation behaves on restriction to that
subgroup. Branching rules can force periods to vanish. For example in the
classical $U(2)$ to $U(1)$ restriction, the representation with highest weight
$(a,b)$, $a\ge b$, contains the weights $b,b+1,\ldots,a$ on
$\operatorname{diag}(z,1)$. This follows by writing it as
$\operatorname{Sym}^{a-b}(\C^2)\otimes\det^b$. A weight $c$ outside that interval
cannot pair with the $c$-component. This illustrates the interlacing conditions
in the note. The proposed chain uses such structure to accumulate multiplicity.

\sources{The underlying suggestions about interlacing, definite groups and length
appear in the \pub{4}{announcement, items 7--10}. The kernel above specifies
the direction of the transform. The figure is a schematic of the proposed construction.}

\sheet{III\enspace The arithmetic construction}{How one lift earns divisibility}{local}
At a finite prime, how much of a lifted form can one divide out? Work in a
discrete valuation ring. Choose a generator $\pi$ of its maximal ideal,
called a uniformizer, so the valuation
is normalized by $v(\pi)=1$, and write $e=v(p)$ for the ramification
index above the rational prime $p$. A vector has normalized valuation at least
$s$ in a lattice $\Lambda$ if it belongs to $\pi^{\lceil es\rceil}\Lambda$.
If a linear map sends $\Lambda$ into $\pi^a\Lambda'$, it sends
$\pi^b\Lambda$ into $\pi^{a+b}\Lambda'$. This addition is how successive
lifts can earn a large common divisor.

\subt{Why one coefficient matters}
The construction uses Hecke operators, which sum a form over arithmetic
translates. At one chosen local component it combines normalized operators
$\widetilde T_i$ into
\[
            Z=\sum_{\ell=0}^{R}(-1)^\ell\beta_\ell\widetilde T_{R-\ell}.
\]
On the lifted vector the positive-index terms have the desired valuation
bound and their coefficients are integral. The zero-index operator is the
identity. For its vector the bound can be smaller by $\alpha$, computed from
weight differences at even positions. Here the coefficient hypothesis is
$v(\beta_R)/e\ge\alpha$. Multiplication restores the missing valuation.
The valuation of a sum is at least the smallest valuation of its terms,
so the entire $Z$-combination meets the stronger bound. This step needs both
the theta-kernel estimates and the coefficient estimate.

For an elementary model, take the vector $(1,1)$ at $p=3$. Suppose the
positive-index term is $9(1,2)$. With identity coefficient $-3$ the sum
is $(6,15)$, divisible by $3$ only once. With coefficient $-9$ it is
$(0,9)$, divisible by $3^2$. The factor required on the identity term can
determine the divisibility of the whole sum.

\subt{Adding the local gains}
Let $J$ index the local components at $p$ and let $V\subseteq J$ be
the components with an auxiliary operator. At stage $j$, write $b_{ij}$
for the basic bound, $g_{ij}$ for the improved bound, $c_{ij}$ for the
cost of moving the operator through the lift, and $d_j$ for the denominator.
On the prescribed input lattice the final normalized valuation is at least
\[
  \sum_{j=1}^{m}
  \left(\sum_{i\in V}\max\{b_{ij},g_{ij}-c_{ij}\}
       +\sum_{i\in J\setminus V}b_{ij}-d_j\right).
\]
To see why the lattice matters, put $s_i\ge\sum_j(g_{ij}-c_{ij}-b_{ij})_+$,
where $z_+=\max(z,0)$. A generator uses an auxiliary operator $m$ times at
each component in a subset of $V$. At each omitted $i\in V$ the input carries
$\lceil es_i\rceil$ extra powers of $\pi$. These cover the possible shortfall
there. The bound then passes to the integral span of the generators.

\sources{The release specifies \src{\growthDir FormalGroupSystem}{352}{the Hecke combination},
\src{\growthDir FormalGroupSystem}{361}{the input generators},
\src{\growthDir FormalGroupSystem}{661}{the one-stage hypotheses} and
\src{\growthDir FormalGroupSystem}{708}{the accumulated lower bound}.
See also the \src{\growthDir MultisetValuation/AofUpdate}{172}{one-step argument}.}

\sheet{III\enspace The arithmetic construction}{Why the periods must share a lattice}{division}
An idele chooses a nonzero element of each completion of a number field,
with a unit at all but finitely many finite places. A Hecke character is a
continuous multiplicative character of the ideles that is trivial on the
diagonally embedded nonzero field elements. It packages local character data
into a global twist. The construction varies this twist while comparing
periods. Rescaling each complex period would destroy the meaning of a
common divisor. One needs a shared integral normalization.

\subt{A division that can be computed}
Take two period columns in the model lattice $\Z^2$ and form
\[
       A=\begin{pmatrix}6&3\\3&6\end{pmatrix}
       =3\begin{pmatrix}2&1\\1&2\end{pmatrix}.
\]
Their determinant is $27$. Division of each column by the common factor $3$
removes $3^2$ from the determinant. The new columns still span an index-$3$
sublattice of $\Z^2$. Thus a residual covolume survives after exact division.
For a column matrix $B$ of full rank, $\det(B^*B)=|\det B|^2$ in the square
case. Pairings of the columns control the volume, and near dependence makes
it small. Both the divisibility and the metric enter a determinant estimate.

\subt{The proposed arithmetic operation}
The note describes a \emph{seesaw} identity. The same integral containing a
theta kernel can be evaluated in two orders on compatible subgroups. In
this application the aim is to identify a Bessel period, a period against
a chosen subgroup and character, shared as the Hecke character varies.
The release makes the normalization precise. For each family member $u$,
the twist $\mu$ must give a period pairing of the form
\[
    P_{u,\mu}(\phi,F)=\zeta_{u,\mu}B_u(\phi,A_\mu F).
\]
Here $B_u$ is a single bilinear form over a number field $K$ and $A_\mu$
maps the test data into its common space. The reference twist has $\zeta=1$.
At each finite place $w$, the scalar for the twist $\mu(w)$ used there is a
unit. Its valuation is zero, so the local divisibility of the period passes
to the common form. The complex comparison also involves isometries and a
nonzero scalar. Its metric information enters the estimate at infinity.

A relative trace formula compares pairings of periods with orbital integrals
to control correlations among columns. A slope inequality of Bost's type compares degrees
of metrized lattices with norms of maps between them. For rank $r$ over
$\Z$, the elementary model of the mean slope is
$-r^{-1}\log(\operatorname{covolume})$. Its sign explains why division and
index losses have to be tracked together.

\sources{See the \pub{4}{announcement, items 12--18},
\src{\aaDir ThetaCorrespondence/ThetaPairLifts/LocalTypeDictionary/SeesawFirstReductionLevel}{25}{the seesaw reduction with local hypotheses}
and \src{Formalism/LocalToGlobal/IndefiniteForms/IndefSide}{383}{the common pairing and Bessel-period conditions}.}

\sheet{III\enspace The arithmetic construction}{Dividing by a fractional ideal}{ideal}
The common divisor has an ideal-theoretic form in the release. Over the
ring of integers $\mathcal O_K$, an ideal need not have a single generator.
Its exponent at each finite place $w$ still makes sense. For one choice $u$
of arithmetic data, the local target accumulated over the lifts is
\[
 s_w(u)=\sum_{j=1}^m\bigl(g_{j,w}+F_{j,w}-E_{j,w}(u)\bigr).
\]
Here $g$ is the usable gain (zero at discarded primes), $F$ the
twisting contribution and $E$ the loss. On the finite support the
integer exponent is $\sigma_w(u)=\lfloor e_w s_w(u)\rfloor$, where
$e_w$ is the ramification index. The floor is taken \emph{after} summation.
Terms from rounding matter once weighted by the logarithmic norms of
many places.

\subt{What is divided and what remains}
For finitely supported integers $\nu_w$ put
$I(\nu)=\prod_w\mathfrak p_w^{\nu_w}$, a fractional ideal of $K$.
The lattice entering the slope comparison is
\[
         N'=I(-\sigma)N_0,\qquad N_0\subseteq I(\tau)M,\qquad
         \operatorname{rank}_K N_0=\operatorname{rank}_K M=r>0 .
\]
With the metrics and totally complex coefficient field of the source,
let $d=[K:\Q]$, $\ell_w=\log N\mathfrak p_w/d$ and
$S=2\sum_w\ell_w(\sigma_w-\tau_w)$. Its slope comparison reads
\[
  \mu_{\max}(N')\le\mu_{\max}(M)+S,\qquad
  \mu(N')=\mu(M)+S-\frac{2}{dr}\log [I(\tau)M:N_0].
\]
Here rank means dimension after extending scalars to $K$.
The symbol $\mu$ denotes the mean slope and $\mu_{\max}$ the maximal slope among
sublattices. The factor $2$ follows the source's metric convention.
The finite index measures what the local divisor has failed to remove.

For example let $K=\Q(i)$, $\pi=2+i$, $M=\Z[i]^2$ and
$N_0=\operatorname{diag}(\pi^3,\pi^2)M$. Take $\sigma=\tau=2$ at
$(\pi)$ and zero elsewhere. Then $N'=\operatorname{diag}(\pi,1)M$,
$S=0$ and the residual index is $5$. In the displayed normalization the
mean slope decreases by $\tfrac12\log5$.

For the actual pairing, each local evaluation must be zero or have valuation
at least $\sigma_w$. An element of $I(-\sigma)$ costs at most $\sigma_w$
powers of the local uniformizer. Their product is integral at $w$ and,
when this holds at every finite place, globally integral. One must still
establish the inclusion, full rank and a uniform bound for the residual index.

\sources{See \src{Formalism/LocalRingCardinality/DataParams\%E2\%82\%82}{212}{the usable gain},
\src{Formalism/LocalRingCardinality/DataDIVDefs}{15}{the exact fractional ideal},
\src{\sizeDir Chain2aCoreNoAnalytic}{128}{the common lattice} and
\src{\sizeDir SlopesHypothesisFinitePlaces}{15}{the slope and index formulas}.}

\sheet{III\enspace The arithmetic construction}{Why the chain has square-root length}{scale}
Consider one real scale $x>1$ in a proposed good family. The data require an
integer $m$ and a small $\e>0$ such that
\[
                 x^{1/2-\e}\ \le m\ \le x^{1/2+\e/2}.
\]
The lower bound is what makes $x/m\le x^{1/2+\e}$. The upper bound limits
terms that grow with $m$. If we put $m=x^\alpha$, the exponents of $m$ and
$x/m$ are $\alpha$ and $1-\alpha$. Their maximum is smallest at $\alpha=1/2$.

\ill{balance}{The solid line is $\max(\alpha,1-\alpha)$, the dashed line the smaller
exponent. This is an exact calculation for two quantities. Additional errors
in the arithmetic comparison are worked out on p.~\pageref{cost}.}

For an example of this balance, take $x=10^6$ and $m=10^3$. Both quantities equal $10^3$.
Shorten the chain to $m=10^2$, and $x/m$ rises to $10^4$. Lengthen it
to $10^4$, and the expense proportional to $m$ is already $10^4$.

\subt{The actual integer in the construction}
Set $\delta_1=x^{-\e/8}$. The prescribed chain length is the smallest even
integer at least $\sqrt{x\delta_1}$, namely
\[
              m=2\left\lceil\tfrac12 x^{1/2-\e/16}\right\rceil .
\]
It exceeds the unrounded power by less than two. For fixed small $\e>0$
the lower bound above follows immediately, and the upper bound follows for
large $x$ since $x^{9\e/16}$ eventually exceeds the rounding factor.
The data, including this integer, may change as $x$ moves through a window.

The announcement reports an exponent $2/3$ in an earlier calculation.
Reaching $1/2+O(\e)$ requires more than choosing a suitable $m$. A fixed
factor of two in an upper bound at \emph{each} step would multiply to
$2^m$, eventually exceeding every fixed power of $x$. The lattice division
on p.~\pageref{ideal} lets one compare the gains and costs at their
actual places. The comparison has to account for the same main terms on
the two sides.

\sources{Compare \src{\sizeDir GrowthBoundGivesRiemannHypothesis}{109}{the admissible range},
\src{Development/AbelianVarietyTheory/Definitions/AutomorphicMembership/EvenCeil}{21}{the rank prescription},
\src{Formalism/LocalRingCardinality/DataParams\%E2\%82\%82}{45}{the scale parameters}
and the \pub{4}{announcement, items 8 and 14--16}.}

\sheet{III\enspace The arithmetic construction}{What cancels in the height estimate}{cancel}
At an ordinary place above $p$, a denominator already occurs in the final
period integral. Let $D_\Gamma$ be a common multiple of the orders of the
finite stabilizers in that sum. The terminal period can lose
$v_p(D_\Gamma)$ in normalized valuation. The column supplied by the lifts
must have valuation at least $s_w(u)+v_p(D_\Gamma)$. Multiplying the two
leaves at least $s_w(u)$ before rounding. The level prime and the two
distinguished type primes have their own local comparisons.

\subt{Comparing the same main term twice}
Write $\mathcal F$ and $\mathcal E$ for the weighted sums of the local
twisting contributions and losses on p.~\pageref{ideal}. The lower bound on
the quotient slope contains $+\mathcal F-\mathcal E$ alongside the
divisibility gain. The upper bound on the dual slope plus the map height
contains the same $+\mathcal F$. Subtraction removes the twisting term.
Finite-place divisibility still leaves a height to estimate at the infinite
places. There one must bound the norms of the composed maps for the same
quotient metrics used to compute the slopes.

For the next cancellation let $A=\sum_t a_t\ge0$ be the total weight,
$\ell_g=\log(g_1g_2)$ for the two distinguished primes, and $T$ the chosen
level. The count of levels involves an average over $f$ weight indices at
each stage. In the source's notation set
\[
 W=\sum_{j=1}^m\frac1f\sum_\rho\log(w_{j\rho}v_0/\pi),\qquad
 D_{\rm cnt}=c_{\rm cnt}L_a(1+m^2/k).
\]
Here $w_{j\rho}$ are weights, $v_0$ is a positive scale, $\pi$ in this
display is the circle constant, and $L_a$ is a logarithmic allowance.
The two required estimates, for every choice over a base index, give
\[
  |\mathcal E-mA\ell_g|\le c_E m B_0,\qquad
  \log T\ge W+m\ell_g-D_{\rm cnt}.
\]
The height contains $\mathcal E-A\log T$. Use the upper bound for
$\mathcal E$ and multiply the lower bound for $\log T$ by $-A$. Then
\[
        \mathcal E-A\log T\ \le\ -AW+c_E m B_0+A D_{\rm cnt}.
\]
The two occurrences of $mA\ell_g$ have exactly the same coefficient.
The term $-AW$ now combines with the archimedean main term in the prime
count. The remaining allowances still depend on $m$ and $k$. Their size
after averaging is the calculation on p.~\pageref{cost}.

\sources{See the \src{Formalism/LocalToGlobal/GroupLattices/AdmissibleAtGeneratorPlaces}{187}{ordinary-place calculation},
\src{\sizeDir QuotientSlopeBoundsWith}{19}{quotient slope, loss and count estimates},
\src{Formalism/RankGapBounds/MeanBoundSecondForm}{118}{dual slope and height},
and \src{\sizeDir ChainSlopeWindows/LevelOfCongruence}{105}{the main-term identity}.
The repository's \href{\repo blob/main/lift-estimate.pdf}{worked lift estimate} follows this calculation from the input lattice to the exponent.}

\sheet{IV\enspace Back to the primes}{Frobenius turns a count into characters}{frobenius}
How does a determinant bound say anything about the $\Lambda$-weighted sum
from p.~\pageref{ramp}? An abelian extension of number fields supplies a
finite abelian group $A$ of symmetries. An unramified rational prime $p$
in an abelian extension of $\Q$ determines an element $\mathrm{Frob}_p$
whose action on the residue field at a prime above $p$ is $z\mapsto z^p$.
In $\Q(i)$ and at odd $p$, this sends $i$ to $i^p$. It is the identity at
$5$, conjugation at $3$, and its square is the identity at $3$. This last
observation accounts for the sign of the $9$ term in our earlier example.

\subt{A prime-power identity}
Let $\widehat A$ be the characters of $A$. For $a,g\in A$, orthogonality gives
\[
        \one_{g=a}=\frac1{|A|}
           \sum_{\xi\in\widehat A}\overline{\xi(a)}\,\xi(g).
\]
When $g=a$ every summand is one. Otherwise multiplying the sum by
$\xi_0(ga^{-1})\ne1$ for a separating character $\xi_0$ permutes its terms,
so it vanishes. Set $g=\mathrm{Frob}_p^r$ and multiply by
$(\log p)(x/p^r-1)$ for $p^r\le x$. A count of prime powers in a prescribed
class becomes a linear combination of character ramps. At ramified primes
there is inertia, the subgroup acting trivially on the residue field.
The release handles those local terms separately in its balance identity.

\subt{Where the fraction $1/(3h+1)$ comes from}
The auxiliary group has a quotient of order $h$. The target proportions on
its cosets are $4/(3h+1)$ at the identity and $3/(3h+1)$ at each other coset.
Equivalently the measure is uniform mass $3/(3h+1)$ per coset plus one
extra unit of mass $1/(3h+1)$ at the identity. Its trivial Fourier coefficient
is $1$. For every nontrivial character the uniform part sums to zero and
the coefficient is $1/(3h+1)$. For $h=2$ this is just $4/7-3/7=1/7$.
In the finite sets used by the release, each observed coset fraction must
differ from its target by at most $x^{-1+\e}$.

In source notation \texttt{ch\_balance} identifies the measured quantity
$\mathrm{ASK}$ minus its explicit main term with a sum of centered character
contributions. The determinant estimates are applied to this difference.
Terms from characters outside the retained family and the errors in the
shares must also be bounded before the answer can be $H_\psi(x)$.

\subt{Why the prime signal occurs $m$ times}
At fixed group indices $u,t$, the comparison replaces each lift's expression
by a common character sum evaluated at $a(t)/b$, within controlled errors.
Call this common term $T(u,t)$. Summing over the $m$ lift labels gives
$\sum_{j=1}^m\re T(u,t)=m\re T(u,t)$. Averaging over the subgroup kills
characters nontrivial on it and supplies the density factor $\delta/2$ on
the next page. This explains the multiplier on the prime sum while allowing
the error to grow with the chain length.

\sources{Compare \src{\laDir ChBalance}{200}{the balance identity and its local cases},
\src{\frDir ExplicitFormulaErrorOfHostConditions}{135}{coset shares},
\src{\prDir TopsSum}{75}{their Fourier coefficients}, \src{\laDir ChCore}{283}{averaging over lifts and group indices}
and \src{\prDir SigALe}{38}{the bound on discarded characters}.}

\sheet{IV\enspace Back to the primes}{Accounting for the four error terms}{cost}
Here is the size calculation after the arithmetic data have supplied the
bound on $\mathrm{ASK}-\mathrm{MAIN}$. Use the release's notation $n$ for
the order of an auxiliary group, $n_T,n_2$ for the sizes of two chosen sets,
$b>0$ for a scale factor, $k$ for a size parameter and $a_E$ for a
nonnegative cost. This $k$ is independent of the derivative order on
p.~\pageref{kernel}. Write $Q_0=n_Tmb$, $\delta=n_2/n$, and let $G$ be
$H_\psi(x)$ with the outside-character contribution $\Sigma$ removed.
The comparison gives
\[
 Q_0\,\frac{\delta}{2}\bigl(G-4x^{1/2+\e}\bigr)
       \le \mathrm{Err}+2Q_0/x.
\]
For $x\ge3$, $\e>0$, $C\ge0$, the size hypotheses include the range of $m$
on p.~\pageref{scale} and
\[
 a_E\le2xb,\quad n_T\ge\tfrac12 n x^{-\e/8},\quad
 x^{-\e}\le\delta\le1,\quad xm\le k\le n^2 x^{1/2-\e}.
\]
With these data $Q_0>0$. The bound assumed for
$\mathrm{Err}/(Q_0x^{C\e})$ is the sum of the four terms in the middle
column. The right column follows by substitution.

\medskip
{\renewcommand{\arraystretch}{1.38}\begin{tabularx}{\linewidth}{@{}l>{\centering\arraybackslash}Xl@{}}
\toprule
\sffamily Source of bound & \sffamily Term after division & \sffamily Upper bound\\
\midrule
$a_E\le2xb$, lower $m$ & $a_E/(mb)$ & $2x^{1/2+\e}$\\
$a_E\le2xb$, $xm\le k$ & $a_E m/(bk)$ & $2$\\
upper $m$, lower $n_T$ & $2(n/n_T)m x^\e$ & $4x^{1/2+13\e/8}$\\
upper $k$, $n_2\le n$ & $(n_2/n_T)(k/n^2)x^{2\e}$ & $2x^{1/2+9\e/8}$\\
\bottomrule
\end{tabularx}}\medskip

For example the third row pays $\e/8$ for $n/n_T$, $\e/2$ for the upper
bound on $m$, and $\e$ explicitly. Their sum is $13\e/8$.
Thus $\mathrm{Err}/Q_0\le10x^{1/2+(C+13/8)\e}$. Dividing the comparison by
$\delta/2$ costs at most $2x^\e$. Since $2/x\le1$, the resulting bound on
$G$ is at most $26x^{1/2+(C+21/8)\e}$. At a sufficiently large scale,
$26\le x^{3\e/8}$, hence $G\le x^{1/2+(C+3)\e}$.

In the cyclic-family comparison one may use $C=e_D+e_L+3$ with
$e_D+e_L\ge0$. The separate estimate is $\Sigma\le4h\,x^{1/2+\e}$.
Once $1+4h\le x^\e$, their sum gives
$H_\psi(x)\le x^{1/2+(e_D+e_L+7)\e}$. The release takes $e_D=2,e_L=0$,
giving the conditional exponent $\tfrac12+9\e$.

\sources{The inputs are \src{\sizeDir GrowthBoundGivesRiemannHypothesis}{109}{SizeWindows},
\src{\laDir ChCore}{283}{the averaged inequality} and
\src{\laDir ChTop}{37}{the pointwise gap estimate}. Compare the
\src{\sizeDir ChainSlopeWindows/LipschitzWeightAndMinimum}{507}{four-term bound} and
\src{\prDir CoreGm}{12}{core\_gm} and
\src{\thirdDir TopShape}{52}{the chosen constants}.}

\sheet{IV\enspace Back to the primes}{A cohomology bound on auxiliary windows}{cohomology}
What could obstruct the division of a family of periods? The announcement
identifies classes left over after twisting. The quantitative input in the
release concerns number fields $E$ chosen inside cyclotomic fields. For a
prime $\ell$ use Galois cohomology unramified away from the primes above
$\ell$, with coefficients $\Q_\ell/\Z_\ell(-j)$. The twist means that Galois
acts by the inverse $j$th power of the cyclotomic character, the character
describing its action on $\ell$-power roots of unity.

In degree one a class can be represented by a cocycle, a function satisfying
$c(gh)=c(g)+g\,c(h)$, taken modulo functions $c(g)=g v-v$. Restriction to
local groups tests whether a global class becomes trivial there. Write
$b_{\ell,j}(E)$ for the $\Z_\ell$-length of the kernel of the release's
localization map at the primes above $\ell$. For a finite module, length
counts powers of $\ell$ in its order. For example
$\Z/\ell^2\Z\oplus\Z/\ell\Z$ has length $3$ and contributes $3\log\ell$
to a logarithmic size estimate.

\begin{statement}
\smallhead{The demanded bound.}\quad Put $t=\lfloor x^{401}\rfloor$.
For the prescribed fields and scales, the lengths for all primes $\ell$
and $1\le j<t$ must be finite and, for \emph{every finite set} $P$ of primes,
\[
      \sum_{\ell\in P}(\log\ell)\sum_{1\le j<t} b_{\ell,j}(E)
                      \le \exp(x^\e/2).
\]
\end{statement}
All terms are nonnegative, so the bound passes to their total $V$. Write
$L=x^\e$, $b=e^L$ and $B_0=fmbL$. The inequality gives $V\le\sqrt b$.
At lift $j$ the congruence correction is bounded by
$d_0[C_0(L B_0+(\log x)^2N_c)+V]$, provided the base rank is $d_0$.
Here $N_c\ge0$ is a second size allowance. Since $(\log x)^2\le L$
eventually and $f,m,L\ge1$, the correction fits
$d_0(|C_0|+1)L(B_0+N_c)$. It must be paid at every lift.

\subt{Which fields, which scales?}
Fix $\psi$, $0<\e<1/52$, and a choice of integers $e(x)$ with
$x^{2+8\e}\le2^{e(x)}\le2x^{2+8\e}$ for $x\ge1$.
The source asks for $J\ge1$ and $X_0$ so that for every $X\ge X_0$ and
every admissible prime pair $q,d$ at scale $X^{1/J}$ there exists $Z$ with
$X^{1/J}\le Z$ and $Z^2\le X$. The bound must hold for every
$x\in[Z,Z^2]$ and every field meeting the prescribed conditions at
$(\psi,e(x),q,d)$. All choices of algebraic closure and local data in the
definition are quantified.

Alpöge calls this a weakening of Vandiver. Classically Vandiver's conjecture
says that a prime $\ell$ does not divide the class number of the maximal
real subfield of $\Q(\zeta_\ell)$ for odd $\ell$. Here the imposed statement
is a length bound for particular twisted groups in varying fields. The release
uses its bound on square windows to supply hosts for the $3/2$-power windows.

\sources{Exact field and prime-pair conditions are in
\src{\chDir ExistsPageN}{12}{VVle, IsChiField, WinPair and VWinAt}.
See \src{Formalism/GaloisCohomology/Statement}{67}{the twist and localization kernel},
\src{Formalism/SharpRestriction/GlobalTheoremInputs}{36}{the per-lift correction},
\src{\chDir DivinNOfHypV}{76}{its absorption} and
\src{\chDir SupplyChi}{8}{the host supply}.}

\sheet{IV\enspace Back to the primes}{Selecting primes and fields for a window}{selection}
The lengths on p.~\pageref{cohomology} depend on the field. Merely choosing
one convenient field would leave the growth in $x$ uncontrolled. The source
makes the choice of auxiliary data quite restrictive.

\subt{The admissible prime pairs}
Write $Y=X^{1/J}$ and $B=Y^{\e/5}$. The primes $q,d$ obey
\[
 \tfrac12 B<\log q\le B,\qquad
 \tfrac12 B(\log Y)^2<\log d\le B(\log Y)^2 .
\]
Further, $q\equiv3\bmod8$, $d\equiv7\bmod16$, $q\mid d+1$, and
$p\mid d+1$ for every odd prime $p\le B$. The cohomological window
statement is conditional on a pair satisfying these constraints. Producing
the hosts used by the prime-sum argument also needs the corresponding
existence and admissibility results.

For the chosen $e=e(x)$, embed $E$ into a cyclotomic field $\Q(\zeta_N)$
with $N$ divisible by $2^{e+2}$, the conductor of $\psi$, $q$ and $d$.
An automorphism $\zeta_N\mapsto\zeta_N^a$ fixes $E$ pointwise precisely when
\[
 a\equiv\pm1\bmod 2^{e+2},\qquad \psi(a)=1,\qquad
 \left(\frac{a}{q}\right)=\left(\frac{a}{d}\right)=1 .
\]
The last two symbols are quadratic characters for the odd primes $q,d$.
This specifies the field used in the host's Frobenius comparisons.
The growing real cyclotomic layers give a concrete example
of a tower of fields in the construction.

\subt{Why a Behrend set enters the picture}
Inside a cyclic group one wants many allowable indices with few close
encounters along certain orbits. The release formulates a Behrend condition
on $T\subset\{0,\ldots,N-1\}$. For nonzero integers $a,b,c$ with
$a+b+c=0$ and $|a|,|b|,|c|\le K$, the only solutions in $T$ to
$a t_1+b t_2+c t_3=0$ must have $t_1=t_2=t_3$. With $K\ge2$ this
excludes a nonconstant three-term arithmetic progression by taking $(1,-2,1)$.
The size requested is at least $N\exp(-C\sqrt{\log N\,\log(K+1)})$ for
an absolute $C>0$.

One concrete use is the source's \texttt{WindowTwo} property. When
$8Lt<n$ for every $t\in T$ and $T$ satisfies this condition at $K=2L$,
a block of $2L$ successive positions in an orbit of additive order $>2L$
in $\Z/n\Z$ hits a fixed translate of $T$ at most twice. A spacing property
and a lower bound on the size of $T$ are both needed. The further transfer
of all these choices into one admissible host family is an arithmetic input.

\sources{See \src{\chDir ExistsPageN}{19}{IsChiField and WinPair},
\src{Formalism/BehrendSets/Statement}{48}{BehrendFree, BehrendBound and WindowTwo},
and the \pub{4}{announcement, items 11, 17 and 19}.}

\sheet{IV\enspace Back to the primes}{Closing the argument at the same scales}{assembly}
For the fixed target $\psi$, write $S_\e$ for the admissible scales of a
chosen auxiliary family. A host supplies the field, characters and Frobenius
data. Its estimates on the retained and discarded characters must refer to
the very same host as the arithmetic chain. Otherwise the prime sum measured
by the determinant comparison could differ from $H_\psi$.

\subt{The order of choices}
The required existence statement supplies constants $e_D,e_L$ and
$\e_1>0$ \emph{before} choosing $\e$. For each $0<\e<\e_1$ it supplies a
compatible family with these two properties.
\begin{enumerate}
\item For every lower cutoff $Y_0$ there is a $Y\ge Y_0$ with
$[Y,Y^{3/2}]\subset S_\e$.
\item There are constants for this family and a threshold $x_0$ such that
the host and chain estimates needed in the comparison apply at every
$x\in S_\e$ above $x_0$. Individual auxiliary objects may vary with $x$.
\end{enumerate}
The comparison of p.~\pageref{cost} supplies a threshold $x_1$ on these
scales. With the release's $e_D=2$ and $e_L=0$, the exponent is
$\tfrac12+9\e$. Given $\theta>\tfrac12$, choose
\[
        0<\e<\min\!\left(\e_1,\frac{\theta-\tfrac12}{9}\right).
\]
Use property 1 with lower cutoff $\max(2,Y_0,x_0,x_1)$. Every point of the
selected window has $H_\psi(x)\le x^{1/2+9\e}\le x^\theta$. This proves
$(\W)$ from the two inputs. If the only available loss exponent grew like
$1/\e$, shrinking $\e$ would not close the exponent gap.

\subt{From the right half of the strip to all characters}
By the criterion of p.~\pageref{separation}, $(\W)$ excludes a zero
$\rho$ of a primitive $\psi$ with $\tfrac12<\re\rho<1$. A primitive zero
with $0<\re\rho<\tfrac12$ gives a zero for $\bar\psi$ at $1-\rho$ by the
functional equation. The gamma and conductor factors are finite and
nonzero in the open strip. Applying the same exclusion to $\bar\psi$ gives
a contradiction. Finally, if $\chi$ modulo $q$ is induced by $\psi$ of
conductor $f$, then
\[
       L(s,\chi)=L(s,\psi)
       \prod_{\substack{p\mid q\\p\nmid f}}(1-\psi(p)p^{-s}).
\]
For $\re s>0$, $|\psi(p)p^{-s}|<1$ in the product. For example the character
induced from $\chi_4$ at modulus $12$ contributes $1+3^{-s}$. This accounts
for every modulus, including $1$, which gives the ordinary Riemann hypothesis.

\sources{See \src{\frDir TopWin}{26}{the window and chain quantifiers},
\src{\chDir ZeroFreeAllOfTopG}{26}{their assembly},
\src{\frDir TopWin}{68}{the choice of $\e$} and
\src{\chDir ClaimStrip}{23}{primitive reduction and reflection}.}

\sheet{V\enspace Reading the release}{How to read the Lean statements}{lean}
Lean represents a proposition as a type. A proof is a term of that type.
The kernel checks the term against the definitions, earlier theorems and
axioms in its environment. Short final declarations can have very long
dependency chains. A useful starting point in this release is a single line
from \texttt{Main.lean}.
\begin{statement}
{\small\ttfamily theorem LFunction\_eq\_zero\_re : DirichletLZerosOnLine\\
\hspace*{1.5em}:= claimAll.2.1}
\end{statement}
The type after the colon names the statement. The expression after the
assignment sign selects a component from the previously proved conjunction
\texttt{claimAll}. To see what that component says, open
\texttt{DirichletLZerosOnLine} in \texttt{Statements.lean}. In mathematical
notation its type is
\[
\begin{gathered}
 \forall N\in\Z_{\ge1},\quad
 \forall\chi:(\Z/N\Z)^\times\to\C^\times\ \text{a character},\quad
 \forall s\in\C,\\[-2pt]
 0<\re s\ \Longrightarrow\ \re s<1\ \Longrightarrow\
 L(s,\chi)=0\ \Longrightarrow\ \re s=\tfrac12 .
\end{gathered}
\]
In Lean, \texttt{[NeZero N]} supplies the nonzero-modulus condition and the
arrows introduce successive hypotheses. The \texttt{LFunction} here is
mathlib's Dirichlet $L$-function. Another declaration,
\texttt{LFunction\_ne\_zero\_of\_half\_lt\_re}, states right-half-plane
nonvanishing with the exception $\chi=\one$, $s=1$. The separate
\texttt{riemannHypothesis} declaration has mathlib's type
\texttt{RiemannHypothesis}.

\subt{Follow one dependency chain}
\begin{enumerate}
\item \src{Supplements/Kernels/MainTheorems}{7}{MainTheorems} obtains
\texttt{allChi} from \texttt{FinalTheorems.everyCharacter} and uses it for
the strip statement. This is where to check the reduction back to every modulus.
\item \src{\chDir MainConclusion}{7}{MainConclusion} supplies both the
arithmetic inputs and \texttt{VWin} using a local--global Galois
compatibility statement called \texttt{Row0b}.
\item \src{\aaDir Foundations/GaloisRepresentations/IntegralStructures/RowSrc}{6}{RowSrc}
obtains that compatibility from a theta-lift statement. Its short proof passes
the problem to imported theorems. A reader auditing this route must examine
the inputs that construct those lifts.
\end{enumerate}
A declaration beginning with \texttt{def} and ending in \texttt{: Prop} can
package a demand as a proposition. Its definition alone gives no witness.
For example, p.~\pageref{cohomology} expands the demand \texttt{VWinAt},
while \texttt{vwin\_of\_row} gives a route to supplying it from another premise.

\sources{Primary entry points include \src{Statements}{3}{Statements.lean},
\src{Main}{4}{Main.lean}, and
\src{\aaDir AutomorphicVarieties/ModularityArgument/ThresholdArguments/DenseThresholdCases/VwinOfRow}{4}{vwin\_of\_row}.
The announcement describes the extended Riemann hypothesis as a future direction.}

\sheet{V\enspace Reading the release}{What has been checked, and what is open}{checks}
\subt{Calculations reproduced here}
The finite-field identity, the path norm bound and the small numerical
examples follow from the sums displayed in the text. On the analytic side,
we derived both centerings, the family residue, the localized tail estimate
for $3/2$-power windows and the contradiction with a pole. On the arithmetic
side, we traced the local coefficient and valuation sum, checked the two
height cancellations and derived $\tfrac12+9\e$ from the size hypotheses.
The final choice of $\e$ and the passage to all moduli follow
once one family supplies the requisite estimates together.

\subt{Questions an arithmetic verification must answer}
\begin{enumerate}
\item Do the chosen theta lifts and test vectors give the required common
period with nonvanishing, compatible integral normalizations and the stated
local factors at every stage? Which inclusions justify the integer exponents
used to divide the period lattice?
\item After division, do the residual lattice indices, correlations between
columns and archimedean norms satisfy the simultaneous slope and error bounds
as $m$ grows? A comparison at one fixed length leaves uniformity open.
\item Does the construction produce hosts meeting the character, cohomology
and chain conditions on each point of the \emph{same} whole window for
arbitrarily small $\e$, with $e_D,e_L$ fixed first?
\end{enumerate}

\subt{What a formal build would check}
The release pins Lean \texttt{v4.33.1} and mathlib revision
\texttt{0df444a360ea}. In the extracted \texttt{lean/} project, the module
target \texttt{lake build +Main} requests the final module and its imports.
After a successful build, a separate file \texttt{Audit.lean} can contain
\begin{statement}
{\small\ttfamily import Main\\
\#print DirichletLZerosOnLine\\
\#print axioms LFunction\_eq\_zero\_re\\
\#print axioms LFunction\_ne\_zero\_of\_half\_lt\_re\\
\#print axioms riemannHypothesis}
\end{statement}
Run \texttt{lake env lean Audit.lean} from that project. The first print
exposes the statement. The axiom queries report transitive dependencies,
including \texttt{sorryAx} if an admitted proof is used. An independent
verification would record the clean build, actual axiom output and a replay
of the compiled proof terms, then compare the formal statements to the
mathematics. This companion reports a source reading and hand calculations.
A full kernel check and an independent derivation of the uniform lift
and cohomological estimates remain open here.

\sources{Links throughout point into the
\href{https://github.com/TomasOrtega/RH/tree/a6ac91a412dcbaab31801355ec5707024a71bf65/extracted/lean}{pinned source mirror}.
See the \href{https://lean-lang.org/doc/reference/latest/ValidatingProofs/}{Lean reference on validation}.
The accompanying README gives the commands for rebuilding every figure and this PDF.}
\end{document}
